% !TEX root = main.tex
\section{Evidence from the Literature and Public Benchmarks}
\label{sec:literature}\label{sec:prisma-results}

\subsection{Search results}
\label{subsec:prisma-flow}

The search identified 150 records, of which 32 external studies or benchmark
sources met the inclusion criteria (Figure~\ref{fig:prisma-flow} and
Table~\ref{tab:prisma-flow}). Most were published in 2024 or 2025, reflecting
the recent introduction of the models considered here. The final extraction
contains 802 rows because benchmark papers report many combinations of task,
model, and metric; the row count should not be interpreted as the number of
independent studies.

\begin{figure}[htbp]
\centering
\includegraphics[width=0.82\textwidth]{prisma_flow.pdf}
\caption{Flow of records through the structured search and extraction.}
\label{fig:prisma-flow}
\end{figure}

\begin{table}[htbp]
\centering
\caption{Structured-search and extraction flow.}
\label{tab:prisma-flow}
\small
\begin{tabularx}{\textwidth}{l r X}
\toprule
Stage & Count & Comment \\
\midrule
Records identified & 150 & Google Scholar, Semantic Scholar, and arXiv \\
Duplicates removed & 30 & DOI or title overlap \\
Titles and abstracts screened & 120 & 50 excluded \\
Full texts assessed & 70 & 38 excluded \\
External studies and benchmark sources included & 32 & Literature extraction \\
Experimental blocks added & 7 & Present study; not literature records \\
Extracted result rows & 802 & Multiple task--model--metric rows per source \\
\bottomrule
\end{tabularx}
\end{table}

Fev-bench and GIFT-Eval account for most rows because their published results
can be separated by task. The extraction includes 20 retail tasks from
fev-bench and four from GIFT-Eval, as well as results for M5, Favorita,
Rohlik, Rossmann, Walmart, and several smaller retail datasets. Foundation
models are consequently over-represented in the raw row count. This imbalance
does not affect the main comparison, which pairs each foundation-model result
with a conventional method in the same task and metric cell.

The three most common measures are \WAPE{}, \MASE{}, and Scaled Quantile Loss.
Their near-equal representation in the paired file---166 comparisons for each
measure---allows metric-specific summaries. Other measures, including the
official M5 \WRMSSE{}, are retained for context but are not pooled with these
three because they answer different forecasting questions.
\label{subsec:descriptive}

\subsection{What the published evidence can compare}
\label{subsec:narrative}

The search revealed a separation between two literatures. Papers introducing
foundation models usually compare them with statistical methods and other
foundation models on broad suites. Retail forecasting papers more often study
LightGBM, XGBoost, or neural models on individual datasets, and many predate
the recent zero-shot models. We found no standalone model paper that reported
a foundation model and a gradient-boosted tree on the same retail series with
the same per-series metric and forecast protocol.

Public benchmark suites partly bridge this gap. Fev-bench and GIFT-Eval report
several model families within common tasks, which makes the 498 descriptive
pairings possible. They also illustrate why the choice of comparator matters:
foundation models tend to have a larger advantage over Seasonal Naive or
Auto~ETS than over LightGBM or CatBoost. Their task-level tables do not,
however, provide the joint prediction errors required to estimate uncertainty
for a difference between models.

M5 provides a further warning against combining results solely by dataset
name. Strong tree models achieve \WRMSSE{} near 0.52--0.54, whereas
published foundation-model results on M5 often use per-series \WAPE{}. The two
sets of numbers cannot be compared directly. A low-volume series with little
test-period demand can dominate mean per-series \WAPE{} but receives a very
different weight under \WRMSSE{}. We therefore treat metric-specific public
results separately and use the matched experiments for the direct retail
comparison.

This division of evidence determines the role of the two empirical analyses
that follow. The local experiments provide tighter comparisons on three
datasets and consumer hardware. The public suites cover a wider range of
tasks, models, and distributional forecasts. Agreement between them is
informative, but neither is taken as a substitute for the other.
\label{subsec:structural-finding}
